\cotefile{10}
\scanpage{17}
\margin{\unsure{\ensuremath{\Vert}}}telle que les formes $\varphi$ sur $M$ qui sont des
$(zP)$ polarisations, sont celles pour lesquelles $\varphi^z$
(défini par $\varphi^z(x,y)=\varphi(x,zy)$) sont \deleted{définies} des
$P$-polarisations.

19) De façon générale, \deleted{\unsure{sont}} pour une $\otimes$-catégorie
$\mathcal M$ sur $K\subset\mathbb R$, si $\varphi$ est une forme
$M\otimes M\unsure{\mathord{-}|\longrightarrow}L_{\unsure{\mathrm{Ligne.}|\mathrm{ligne.}}}$, \deleted{\unsure{alors}} et $z\in {}_2\mathrm{Aut}_{\otimes}(\mathrm{id}_{\mathcal M})={}_2Z(K)$
($Z=$ Centre du groupe de \unsure{Galois|choix} \deleted{\unsure{quelconque}} de $\mathcal M$),
alors on définit $\varphi^z$ sur $M\times M$ par
$\varphi^z(x,y)=\varphi(x,zy)$. On trouve alors, si $\varphi^z$ et $\psi$ (sur $N$)
sont non dég. et $\epsilon$-sym ($\epsilon=\pm1$), et si
% slanted
\margin{La condition $z^2=1$ sert \unsure{à|:} assurer que $\varphi$ $\epsilon$-sym implique $\varphi^z$ $\epsilon$-sym.}
\deleted{\unsure{g}} $u:M\longrightarrow N$, $u^{*(\varphi,\psi)}:N\longrightarrow M$, alors
$u^{*(\varphi^z,\psi^z)}=z_M^{-1}u^{*(\varphi,\psi)}z_N=u^{*\varphi,\psi}$. Donc : \deleted{$\varphi$ et $\psi$ sont}

\[
\begin{aligned}
\varphi\ \text{de Weil}&\Longleftrightarrow\varphi^z\ \text{de Weil}\\
\varphi\ \text{et}\ \psi\ \text{de Weil comp.}&\Longleftrightarrow\text{idem pour }\varphi^{\unsure{\epsilon|z}}\text{ et }\psi^\epsilon.
\end{aligned}
\]

D'ailleurs on a $\varphi^\epsilon\otimes\psi^\epsilon=(\varphi\otimes\psi)^\epsilon$. \deleted{\unsure{Donc on}}
\deleted{On en conclut}
\deleted{bien que} Si $z$ opère par id sur $K(1)$
(i.e. $z\in Z'(K)$) alors on voit que si $P$ est
une polarisation, alors les $z(P(M))=\{z\varphi\mid\varphi\in P(M)\}$
($M\in\mathrm{Ob}\,\mathcal M^i$; $i\in\mathbb Z$) forment une polarisation $\deleted{\unsure{z}}\unsure{\epsilon|z}P$ sur $\mathcal M$.
On voit tout de suite que \deleted{cette} dernière
n'est égale à $P$ que si $z=1$. Ainsi,
\margin{? \unsure{!!|\ensuremath{\Vert}}}
$ {}_2Z(\mathbb R)$ \uline{opère librement sur} $\mathrm{Pol}(\mathcal M)$. \uuline{Question}
$\mathrm{Pol}(\mathcal M)$ \unsure{est-il|est il} un pseudotorseur sous $ {}_2Z(\mathbb R)$ ?
% archive sheet number: 16
